\section{Veicoli veloci: sensori a lunga portata, droni e aerei}\label{sec:lp}

I capitoli precedenti descrivono il sistema con una camera di profondità ferma, fino a $3\unita{m}$. Questo capitolo porta gli stessi principi a un'altra scala: oggetti che si muovono da qualche metro al secondo fino a $210\unita{m/s}$, osservati da sensori a lunga portata. Il caso centrale è un \emph{sito aerospaziale sorvegliato}: un piazzale, un'area di prova o il varco di un hangar, con droni in volo e aerei in rullaggio. Gli aerei in volo sono trattati come limite dichiarato (\S\ref{sec:lp-volo}).

Tutto ciò che segue è calcolo o simulazione. I sensori sono simulati dai valori delle schede tecniche; il moto viene da scenari scritti a mano o da traiettorie pubbliche reali, con rumore di misura e occlusioni simulati. Nessuna prova usa un sensore fisico, e nessuna rilevazione reale di sensore è stata elaborata. Le etichette di classe sono assegnate a mano: il classificatore resta un componente richiesto e non realizzato (\S\ref{sec:sti}). Il sistema osserva e avvisa un operatore: nessuna uscita è un comando, e non vi sono funzioni d'ingaggio, di puntamento o di contrasto.

I principi sono gli stessi dei capitoli~\ref{sec:map} e~\ref{sec:occ}: ciò che non è osservato non è mai «libero»; l'avviso di possibile presenza si toglie solo guardando, o rivedendo l'oggetto; nessuna etichetta di classe riduce un avviso. Cambia ciò che serve per tenerli veri a questa scala, e il capitolo mostra anche dove, in simulazione, non lo sono.

\subsection{Il problema di scala}\label{sec:lp-scala}

Un avviso serve se arriva prima che chi lo riceve abbia bisogno di agire. Per un ricevente fermo e un oggetto che si avvicina con velocità relativa $v_{\mathrm{rel}}$, la distanza a cui l'avviso deve partire è
\begin{equation}\label{eq:lp-davv}
d_{\mathrm{avv}} = v_{\mathrm{rel}}\,\bigl(\tau_{\mathrm{lat}}+T_r+T_{\mathrm{man}}\bigr)+d_{\mathrm{man}}+m ,
\end{equation}
con $\tau_{\mathrm{lat}}$ la latenza del sistema, $T_r$ il tempo di reazione del ricevente, $T_{\mathrm{man}}$ il tempo della manovra o della messa in sicurezza, $d_{\mathrm{man}}$ lo spazio di manovra (per esempio la frenata $v^2/2a_f$) e $m$ un margine. Per un ricevente in moto vale la forma della \eqref{eq:lim-dserve}, con la velocità massima dell'inviluppo al posto di quella del soggetto. L'avviso è \emph{tempestivo} se la traccia è confermata, o la zona è dichiarata «libera», a una distanza non minore di $d_{\mathrm{avv}}$; altrimenti l'uscita dice «avviso non tempestivo per $\langle$classe$\rangle$».

\begin{table}[H]\centering\footnotesize
\caption{Distanza d'avviso richiesta in sei scenari nominali, dalla \eqref{eq:lp-davv} \prova{C}. Latenze, tempi di reazione e margini sono in gran parte ipotesi di progetto; il tempo di reazione su strada è preso da una fonte secondaria \cite{wiki-stopping-sight}. Nel primo scenario il $78\,\%$ della distanza viene dai $30\unita{s}$ ipotizzati per mettere in sicurezza l'area.}\label{tab:lp-scala}
\begin{tabularx}{\linewidth}{@{}>{\RaggedRight\arraybackslash}p{3.3cm}>{\RaggedRight\arraybackslash}p{3.2cm}Lp{1.3cm}>{\RaggedRight\arraybackslash}p{1.9cm}@{}}
\toprule
Scenario & Ricevente & Tempi e margini & $v_{\mathrm{rel}}$ & $d_{\mathrm{avv}}$\\
\midrule
drone (15--$21\unita{m/s}$) dietro un edificio del sito & operatore di sito & latenza $1\unita{s}$, reazione $5\unita{s}$, $30\unita{s}$ di messa in sicurezza, margine $50\unita{m}$ (ipotesi) & $21\unita{m/s}$ & $806\unita{m}$\\
aereo in rullaggio ($9\unita{m/s}$) dietro un hangar & conducente di un mezzo che attraversa la via & $1\unita{s}$ + $2{,}5\unita{s}$; frenate $1{,}0$ e $1{,}5\unita{m/s^2}$; margine $23\unita{m}$ & $17{,}3\unita{m/s}$ & $147\unita{m}$; tramite la torre $277\unita{m}$\\
aereo in finale e drone sotto la copertura radar & operatore di torre & $4+10+12{,}5\unita{s}$, più $10\unita{s}$ di riattaccata; margine $610\unita{m}$ & $91\unita{m/s}$ & $3{,}9\unita{km}$\\
drone con i sensori a bordo di un altro drone (solo analisi) & pilota remoto; sistema di bordo & $0{,}5+3\unita{s}$; di bordo $0{,}4\unita{s}$ & $21\unita{m/s}$ & $126\unita{m}$; $60\unita{m}$\\
auto a un incrocio dietro un camion & conducente; sistema di bordo & $0{,}1+2{,}5\unita{s}$, $3{,}4\unita{m/s^2}$; di bordo $0{,}5\unita{s}$, $6{,}9\unita{m/s^2}$ & $13{,}9\unita{m/s}$ & $67\unita{m}$; $24\unita{m}$\\
treno in stazione dietro un convoglio fermo; treno in linea & personale e macchinista & $1+3\unita{s}$, frenata $0{,}8\unita{m/s^2}$; in linea $600\unita{m}$ di frenata d'emergenza & $16{,}7$; $44{,}4\unita{m/s}$ & $261\unita{m}$; $798\unita{m}$\\
\bottomrule
\end{tabularx}
\end{table}

Dalla tabella~\ref{tab:lp-scala} seguono due fatti \prova{C}. Il primo: la portata di $3\unita{m}$ non entra in nessuno scenario, perché la distanza minima richiesta è $24\unita{m}$; serve un altro ordine di grandezza, da $10^2\unita{m}$ sul piazzale a $10^3\unita{m}$ e oltre in avvicinamento. Il secondo: nelle geometrie scelte l'occlusione cade dentro la distanza d'avviso. Il drone riappare a $375\unita{m}$ dall'operatore, che ne chiede $806$; l'aereo esce dall'ombra dell'hangar a $63\unita{m}$ dall'incrocio, contro $147$ richiesti. Un sistema che avvisa solo quando rivede l'oggetto arriva tardi: l'avviso deve partire prima dell'occlusione e restare acceso per tutta la sua durata. Il secondo fatto vale per le geometrie della tabella, non in generale.

\subsection{Sensori simulati dalle schede}\label{sec:lp-sensori}

La tabella~\ref{tab:lp-sensori} riassume i sensori simulati. I valori vengono dalle schede dei costruttori \cite{robin-iris,terma-scanter5000,hesai-at128,livox-tele15}, dalle norme statunitensi sull'ADS-B e sull'identificazione remota dei droni \cite{cfr-91-227,cfr-89-310} e da conti; ogni campo senza fonte è un'ipotesi. In tutto i parametri sono 124, di cui 57 ($46\,\%$) ipotesi: tutte le probabilità di rilevamento, i tassi di falso allarme, le latenze e i costi. Le schede non dichiarano a quale probabilità di rilevamento si riferisce la portata.

\begin{table}[H]\centering\footnotesize
\caption{Sensori simulati e valori usati. (i): ipotesi, senza una fonte aperta. Le portate valgono per la sezione o la riflettività indicate dalla scheda, non per qualunque oggetto.}\label{tab:lp-sensori}
\begin{tabularx}{\linewidth}{@{}>{\RaggedRight\arraybackslash}p{3.3cm}>{\RaggedRight\arraybackslash}p{3.0cm}>{\RaggedRight\arraybackslash}p{2.0cm}L>{\RaggedRight\arraybackslash}p{1.6cm}@{}}
\toprule
Sensore (esemplare) & Portata & Campo & Errore di misura & Rivisita\\
\midrule
radar per droni (Robin IRIS) \cite{robin-iris} & $1{,}1$--$7{,}6\unita{km}$ secondo il drone; strumentata 5 o $12\unita{km}$ & $360^\circ\times60^\circ$ & $0{,}6\unita{m}$; $0{,}6$--$0{,}75^\circ$; $1^\circ$ in elevazione & $1\unita{s}$\\
radar a pannello (tipo EchoGuard) & (i) & (i) & $<3{,}25\unita{m}$; $<1^\circ$; $<1{,}5^\circ$ in elevazione \cite{echodyne-echoguard} & (i)\\
radar del movimento a terra (Terma SCANTER 5000) \cite{terma-scanter5000} & $5\unita{km}$ per $1\unita{m^2}$ (i); $9{,}6\unita{km}$ in avvicinamento & $360^\circ$, senza quota & cella di $3\unita{m}$; $0{,}29^\circ$ (conto) & $1\unita{s}$\\
LiDAR a campo largo (Hesai AT128) \cite{hesai-at128} & $210\unita{m}$ al $10\,\%$ di riflettività & $120^\circ\times25{,}4^\circ$ & $0{,}1^\circ\times0{,}2^\circ$; $3\unita{cm}$ (i) & $0{,}1\unita{s}$ (i)\\
LiDAR a campo stretto (Livox Tele-15) \cite{livox-tele15} & $320\unita{m}$ al $10\,\%$, $500\unita{m}$ al $50\,\%$ & $14{,}5^\circ\times16{,}2^\circ$ & $2\unita{cm}$; $<0{,}03^\circ$ & $0{,}1\unita{s}$\\
radar 4D per veicoli & $300\unita{m}$ (i: scheda non aperta) & $120^\circ$ (i) & $0{,}22\unita{m}$; $1{,}2^\circ$; $0{,}1\unita{m/s}$ (i) & $0{,}05\unita{s}$ (i)\\
ADS-B \cite{cfr-91-227} & orizzonte del ricevitore & solo chi trasmette & $<92{,}6\unita{m}$; $<10\unita{m/s}$ & $\le1\unita{s}$; latenza $\le2\unita{s}$\\
identificazione remota dei droni \cite{cfr-89-310} & circa $1\unita{km}$ (i) & solo chi trasmette & $30{,}5\unita{m}$ al $95\,\%$; quota $45{,}7\unita{m}$ & $\le1\unita{s}$; latenza $\le1\unita{s}$\\
camera nel visibile, 12 Mpixel, $300\unita{mm}$ (i) & drone di $0{,}3\unita{m}$ a $13\unita{km}$, limite geometrico \cite{wiki-johnson} & $2{,}7^\circ$ & $11{,}5\unita{\mu rad}$ per pixel; nessuna distanza & $0{,}1\unita{s}$\\
camera termica $640\times480$, $100\unita{mm}$ (i) & drone a $1{,}2\unita{km}$, limite geometrico & $4{,}4^\circ$ & $120\unita{\mu rad}$; nessuna distanza & $0{,}1\unita{s}$\\
\bottomrule
\end{tabularx}
\end{table}

\paragraph{Portata del radar e sezione radar.} La scheda del radar per droni dà portate per classe di drone, senza la sezione radar. Attribuendo $0{,}01\unita{m^2}$ (i) al drone da $0{,}9\unita{kg}$ visto a $3100\unita{m}$, la portata scala come $R(\sigma)=3100\unita{m}\cdot(\sigma/0{,}01\unita{m^2})^{1/4}$: $3{,}1\unita{km}$ per un multirotore piccolo, $9{,}8\unita{km}$ per un velivolo leggero di $1\unita{m^2}$, $12\unita{km}$ (limite strumentato) per un aereo di linea \prova{C}. Un drone da $0{,}249\unita{kg}$ si vede solo entro $1{,}1\unita{km}$.

\paragraph{Probabilità di rilevamento (modello, i).} Per un oggetto fluttuante si usa il modello di Swerling~1 per un impulso, $P_d=P_{fa}^{1/(1+\mathrm{SNR})}$, con $\mathrm{SNR}\propto\sigma/R^4$. Con $P_d=0{,}5$ alla portata dichiarata e $P_{fa}=10^{-6}$ il rapporto segnale-rumore di riferimento è $18{,}9$ ($12{,}8\unita{dB}$). Se la portata di scheda valesse invece a $P_d=0{,}9$, la portata utile sarebbe il $62\,\%$ di quella usata: è l'ipotesi che pesa di più sull'intero modello.

\paragraph{Che cosa vuol dire «non osservato» a questa scala.} Una zona è non osservata per cinque cause, che l'uscita distingue: fuori campo; senza linea di vista (dietro un hangar, sotto l'orizzonte radar); \emph{sotto la soglia}, cioè in campo e in vista ma con una probabilità di rilevamento insufficiente per \emph{quella} classe; fra due rivisite (a $70\unita{m/s}$ un aereo si sposta di $70\unita{m}$ in un secondo, più di cento volte l'errore in distanza); oggetto che non coopera o non emette. La terza causa è nuova rispetto alla camera di profondità: la stessa zona può essere osservata per un aereo e non osservata per un drone piccolo. I sensori cooperativi (ADS-B, identificazione remota) aggiungono presenze e identità, ma il loro silenzio non osserva niente: un trasponditore spento non è un'assenza.

\subsection{Rappresentazione per sensore e «libero» per classe}\label{sec:lp-libero}

Ogni sensore ha una propria griglia polare in distanza, azimut ed elevazione, con i quattro stati del capitolo~\ref{sec:map}: libero, occupato, in ombra, non osservato. L'ombra è calcolata dagli ostacoli in pianta e dalla maschera del terreno. Il «libero» non si scrive più da un solo fotogramma, perché un radar che ruota illumina una cella una volta per giro e con una probabilità di rilevamento minore di uno: si scrive per classe, cumulando le passate.

\begin{definizione}[«Libero» per classe]\label{def:lp-libero}
Per una classe $c$ di sezione $\sigma_c$ e velocità massima $v_{\max,c}$, sia
\[W_c=\min\bigl(t_{\max},\,\ell/v_{\max,c}\bigr)\]
la finestra oltre la quale un oggetto può aver cambiato cella di lato $\ell$, e sia $n$ il numero di passate che hanno illuminato la cella in $[t-W_c,t]$ dopo l'ultimo ritorno. La probabilità cumulata di averlo visto è
\[
P_{c}=1-\bigl(1-P_d(R,\sigma_c)\bigr)^{n_{\mathrm{eff}}},\qquad n_{\mathrm{eff}}=1+(n-1)(1-\rho_c),
\]
con $\rho_c$ la correlazione fra passate, posta a 1 in mancanza di una misura ($n_{\mathrm{eff}}=1$), e presa come minimo sulla cella e sulle otto vicine. La cella è «libera per $c$» se non è in ombra, la linea di vista è accertata e $P_c\ge0{,}99$. È «libera» se lo è per ogni classe ammessa dallo scenario; altrimenti l'uscita dice «zona non verificata per $\langle$classe$\rangle$».
\end{definizione}

\begin{proposizione}[La regola con $n_{\mathrm{eff}}=1$ è prudente per ogni correlazione]\label{prop:lp-correlate}
Siano $D_1,\dots,D_n$ gli eventi di rilevamento nelle $n$ passate, con la stessa probabilità marginale $P_d$ e dipendenza qualunque. Allora
\[\Pr(\text{nessun rilevamento})\le 1-P_d :\]
la probabilità di non vedere un oggetto presente non supera quella che la regola con $n_{\mathrm{eff}}=1$ dichiara.
\end{proposizione}
\begin{proof}
L'evento «nessun rilevamento nelle $n$ passate» è contenuto nell'evento «nessun rilevamento nella prima», che ha probabilità $1-P_d$. Nella forma del modello a sezione fluttuante, con $P_d$ condizionata a $\sigma$ variabile, vale $\mathbb{E}[(1-P_d)^n]\le\mathbb{E}[1-P_d]$ perché $0\le1-P_d\le1$.
\end{proof}
La proposizione è \prova{D} nel modello delle passate. Il prezzo è la portata: per un multirotore di $0{,}01\unita{m^2}$, con un radar di riferimento di $P_d=0{,}9$ a $5\unita{km}$ su $1\unita{m^2}$, rivisita di $1\unita{s}$ e celle di $100\unita{m}$, il «libero» arriva a $2360\unita{m}$ se si contano 5 passate indipendenti e a $877\unita{m}$ con $n_{\mathrm{eff}}=1$ \prova{C}. Contare le passate come indipendenti è sbagliato quando l'aspetto dell'oggetto non cambia, per esempio un drone fermo in volo: con la sezione tenuta per 5 passate la probabilità di mancarle tutte alla portata di $2360\unita{m}$ è $0{,}263$, contro $0{,}010$ dichiarato (Monte Carlo di $400\,000$ casi) \prova{S}.

\paragraph{Più sensori.} Lo stato fuso segue l'ordine occupato, libero, in ombra, non osservato. Fra sensori di tipo diverso la probabilità di mancare l'oggetto è il prodotto delle singole; fra sensori dello stesso tipo si prende la migliore, perché l'aspetto può essere correlato (ipotesi). Un sensore con filtro Doppler non produce mai «libero» e il suo silenzio non conta per un oggetto con velocità radiale sotto la soglia del filtro; ogni ritorno toglie il «libero» alle celle entro $3\sigma$ dal punto misurato, più una cella. Un sensore in perdita di affidabilità rende tutta la sua griglia «non osservata». Un ritorno recente di un radar vieta il «libero» attorno al punto misurato per $5\unita{s}$: una sola passata mancata non basta ad aprirlo. Il prezzo di questa ritenuta, cioè le chiusure tardive dell'avviso, non è misurato.

\paragraph{La mappa come ingresso di sicurezza.} Una passata conta come sguardo solo se la linea di vista è accertata: o la mappa degli ostacoli è dichiarata completa, o sul raggio era atteso un ritorno dello sfondo oltre la cella ed è stato misurato alla distanza attesa entro $2\unita{m}$. Se lo sfondo manca il raggio non è accertato; se è più vicino dell'atteso c'è un ostacolo non in mappa, e le celle oltre vanno in ombra. In mancanza di un'indicazione esplicita la mappa non è dichiarata completa, e non c'è nessun «libero». Il cielo non ha sfondo: lì vale solo la mappa dichiarata. Su casi costruiti la regola regge: con un edificio assente dalla mappa e la mappa non dichiarata completa la cella dietro resta non osservata, e con lo sfondo misurato a $400$ invece che a $1500\unita{m}$ va in ombra e l'ostacolo è annotato \prova{S}. Se invece la mappa è dichiarata completa ed è sbagliata, un edificio che non vi compare apre un «libero» falso: è un controesempio (5 casi costruiti su 6). Per questo nel prototipo, per impostazione, la mappa non è dichiarata completa, e nei sei casi la cella dietro l'edificio resta «non osservata» \prova{S}. La convalida della mappa con lo sfondo misurato non è ancora usata nel simulatore (\S\ref{sec:lp-smentiti}).

\paragraph{Classi escluse.} Lo scenario dichiara le classi ammesse e quelle escluse; per le escluse l'uscita dice sempre «zona non verificata». Nel sito sorvegliato è escluso il drone sotto i $250\unita{g}$, che il radar non vede oltre circa $1{,}1\unita{km}$, e sono esclusi gli aerei in volo. Il «libero», quindi, vale solo per le classi dichiarate, con la mappa data e con passate trattate come correlate.

\subsection{Inseguimento multisensore}\label{sec:lp-inseguimento}

Lo stato di una traccia è posizione e velocità in tre dimensioni nel riferimento fisso Est--Nord--Alto. Il filtro è una miscela di modelli interagenti \cite{barshalom2001} (velocità quasi costante, virata coordinata a sinistra e a destra, manovra), con parametri per classe; una traccia senza classe usa tutti i modelli con i parametri dell'inviluppo. Le misure di tutti i sensori si fondono una alla volta, nell'ordine d'arrivo: previsione all'istante della misura, porta sulla distanza di Mahalanobis della posizione ($d^2\le16{,}27$, quantile $0{,}999$ di un chi quadro con 3 gradi di libertà), assegnazione, aggiornamento in forma di Joseph. Una misura appena fuori porta ($16{,}27<d^2\le30{,}66$) non apre una traccia: diventa una «misura dubbia», e solo la terza in $5\unita{s}$ apre una traccia nuova. La conferma chiede 2 misure su 3. Le misure cooperative portano un identificativo: lo stesso identificativo è la stessa traccia, identificativi diversi non si associano mai, e una traccia vista solo dal sensore cooperativo è dichiarata «non verificata da un sensore». Una traccia si chiude solo come «non più rilevata»: in vista, con la probabilità cumulata dei mancati rilevamenti, per la classe più piccola ammessa, almeno $0{,}99$. In ombra, fuori copertura o con il sensore non affidabile non si accumula nulla, e la traccia resta con l'avviso.

In simulazione \prova{S}, con un aereo a $70\unita{m/s}$ in virata per $15\unita{s}$, un radar a $1\unita{Hz}$, una camera a $5\unita{Hz}$ e un sensore cooperativo a $2\unita{Hz}$ spento a metà prova, l'errore quadratico medio di posizione è $39{,}9\unita{m}$ con il solo radar e $4{,}4\unita{m}$ con la fusione, che è migliore in 40 semi su 40, con 0 cambi d'identità. Con due oggetti senza classe che si incrociano a $100\unita{m}$ di quota di separazione, gli scambi d'identità sono 0 su 40 semi $[0;\,0{,}088]$. Sono casi costruiti senza falsi ritorni: con il disturbo e con più di due oggetti l'identità non è provata.

\subsection{Insieme di possibile presenza e spegnimento dell'avviso}\label{sec:lp-presenza}

L'avviso è l'insieme di possibile presenza del capitolo~\ref{sec:occ}, calcolato con l'\emph{inviluppo dello scenario}: la velocità e l'accelerazione massime fra tutte le classi ammesse, mai quelle dell'etichetta. Con centro nell'ultima posizione stimata, $\tau$ il tempo dall'ultima misura e $v_0=\min(|\hat v|+3\sigma_v,\,v_{\max})$, il raggio è
\begin{equation}\label{eq:lp-raggio}
r(\tau)=R(\tau;v_0,v_{\max},a_{\max})+3\sigma_p+w,\qquad
R=\begin{cases} v_0\tau+\tfrac12 a_{\max}\tau^2 & \tau\le\tau_1,\\ R(\tau_1)+v_{\max}(\tau-\tau_1) & \tau>\tau_1,\end{cases}
\end{equation}
con $\tau_1=(v_{\max}-v_0)/a_{\max}$ e $w=50\unita{m}$ (ipotesi). L'insieme è una sfera, salvo che lo scenario ammetta solo classi vincolate a un binario: allora è un tratto di binario, che a $30\unita{s}$ riduce una sfera di $2{,}5\unita{km}$ di raggio a un tratto di $675\unita{m}$ per il treno ad alta velocità. Il vincolo vale perché il treno non può lasciare il binario, non per l'etichetta.

\begin{proposizione}[Contenimento nell'inviluppo]\label{prop:lp-contenimento}
Se l'oggetto ha velocità al più $v_{\max}$ e accelerazione al più $a_{\max}$, e la sua posizione e velocità all'ultima misura stanno entro $3\sigma_p$ e $3\sigma_v$ dalla stima, allora a ogni $\tau\ge0$ l'oggetto sta nella sfera di raggio $r(\tau)$ della \eqref{eq:lp-raggio}.
\end{proposizione}
\begin{proof}
La velocità al tempo $s$ è al più $\min(v_0+a_{\max}s,\,v_{\max})$; integrando fra $0$ e $\tau$ si ottiene $R(\tau)$, a cui si aggiunge l'errore di posizione iniziale.
\end{proof}
La proposizione è \prova{D} nel modello e vale quanto valgono $v_{\max}$ e $a_{\max}$, che sono ipotesi o valori presi dai dati. Su $60\,000$ traiettorie simulate entro i limiti dichiarati (10 classi, 2000 traiettorie, 3 orizzonti) nessuna è uscita dall'insieme \prova{S}; è un controllo del conto, non una frequenza. L'inviluppo usato nel sito sorvegliato comprende i multirotori, i droni ad ala fissa e gli aerei in rullaggio; per i multirotori veloci assume $45\unita{m/s}$ e $25\unita{m/s^2}$, valori di ipotesi senza fonte. Per gli aerei in volo l'inviluppo è portato a $210\unita{m/s}$, perché i dati pubblici mostrano decolli a più di $200\unita{m/s}$ (\S\ref{sec:lp-traiettorie}). Su casi costruiti, un drone che scatta da fermo con $20$--$25\unita{m/s^2}$ fino a $45\unita{m/s}$ e un aereo a $205\unita{m/s}$ restano dentro l'avviso a ogni $\tau$ \prova{S}; oltre l'ipotesi, a $60\unita{m/s}$, il drone esce dall'avviso dopo $8\unita{s}$.

\paragraph{Spegnimento prudente.} L'avviso si spegne solo in due casi: l'oggetto è rivisto, oppure ogni cella che tocca l'insieme di possibile presenza è «libera» per ogni classe ammessa, con passate successive alla perdita. Il controllo è esatto sulle celle: una griglia affidabile deve contenere tutta la sfera, e ogni cella con il centro entro il raggio più mezza diagonale deve essere libera. Non c'è alcuno spegnimento a tempo. Poiché l'inviluppo non dipende dalla classe, nessuna etichetta riduce l'avviso, per costruzione \prova{D}; nel prototipo, su un passaggio, l'avviso è identico con e senza etichette in 228 cicli su 228, con l'inviluppo precedente degli aerei (cicli correlati: è un controllo, non una frequenza) \prova{S}.

\paragraph{Perché la regola ingenua cade.} Una regola più semplice conta come prova d'assenza una passata senza ritorno sul punto previsto dell'oggetto, o su un campione di punti dell'insieme, e chiude la traccia dopo abbastanza passate. Ma l'oggetto nascosto non è nel punto previsto: è da qualche parte nell'insieme, che a $10\unita{s}$ ha per un multirotore un raggio di circa $200\unita{m}$. In simulazione, con quella regola, l'avviso si spegne senza che si sia guardata tutta la zona in 564 passaggi su 2400 $[0{,}218;\,0{,}252]$; con la regola prudente, sugli stessi semi, in 0 su 600 $[0;\,0{,}006]$ \prova{S}. Con la regola prudente nel sistema, 0 passaggi su 200 per configurazione restano senza regione.

\paragraph{La zona prevista.} Oltre all'avviso, per una traccia con classe il sistema emette una zona di previsione al livello $0{,}90$: una sfera centrata sulla previsione a velocità costante, con raggio dato da un quantile conforme dell'errore \cite{angelopoulos2021}, tarato per classe a 1, 2, 5, 10 e $30\unita{s}$, più $3(\sigma_p+\sigma_v\tau)$. La zona usa l'etichetta, e per questo non è l'avviso. Diventa «previsto ampio» oltre l'ultimo orizzonte tarato, quando il suo raggio supera $0{,}8$ di quello dell'avviso, per una traccia senza classe o con un sensore non affidabile. Per un multirotore l'insieme raggiungibile ha raggio $105\unita{m}$ a $5\unita{s}$ e $210\unita{m}$ a $10\unita{s}$ \prova{C}: oltre pochi secondi la zona non dice nulla di più dell'avviso, e l'avviso si toglie solo tornando a guardare il volume, con un altro sensore o da un altro punto di vista.

\subsection{Previsione su traiettorie pubbliche reali}\label{sec:lp-traiettorie}

Il moto reale viene da sette archivi pubblici (tabella~\ref{tab:lp-dati}). Il rumore di misura è simulato dalle schede, le occlusioni sono inserite a mano (10 istanti a caso per tratto e per durata), e c'è un solo oggetto per scena, senza associazione. Su metà dei tratti si tarano il rumore di processo e il raggio della zona; le cifre sono sull'altra metà. Poiché più tratti appartengono allo stesso mezzo, la divisione è stata rifatta \emph{per mezzo}, e gli intervalli sono un bootstrap a grappoli per mezzo (2000 ricampionamenti); un Wilson con $n$ uguale al numero di mezzi dà un limite ancora più prudente. La verità è la posizione di bordo (GNSS o ADS-B), non una misura esterna.

\begin{table}[H]\centering\footnotesize
\caption{Archivi pubblici di traiettorie usati per il moto. I tratti sono quelli convertiti; il passo è quello della fonte.}\label{tab:lp-dati}
\begin{tabularx}{\linewidth}{@{}>{\RaggedRight\arraybackslash}p{2.2cm}L>{\RaggedRight\arraybackslash}p{1.6cm}>{\RaggedRight\arraybackslash}p{1.6cm}>{\RaggedRight\arraybackslash}p{3.3cm}@{}}
\toprule
Classe & Fonte e licenza & Tratti & Passo & Verità\\
\midrule
aerei & ricezione ADS-B attorno ad Atlanta, 13 minuti \cite{adsblol}; ODbL 1.0 & 53 & 3,1--$5{,}1\unita{s}$ & ADS-B, non elaborato\\
aerei di linea & traiettorie ADS-B Londra--Zurigo 2025 \cite{zenodo-lhr-zrh}; CC BY 4.0 & 80 & $1\unita{s}$ & ADS-B filtrato e interpolato dalla fonte\\
aviazione generale & TrajAir \cite{trajair}; CC BY 4.0 & 60 & $1\unita{s}$ & ADS-B elaborato dalla fonte\\
multirotore & voli di un quadricottero per consegne \cite{cmu-drone}; CC BY 4.0 & 196 & $0{,}12\unita{s}$ & GNSS e inerziale di bordo\\
ala fissa senza pilota & ALFA \cite{alfa}; CC BY 4.0 & 47 & $0{,}24\unita{s}$ & GNSS dell'autopilota; guasti indotti\\
auto & NGSIM, US-101 \cite{ngsim}; CC BY-SA 3.0 & 157 & $0{,}1\unita{s}$ & telecamere dall'alto\\
treni & posizioni dei treni finlandesi, 12 novembre 2026 \cite{digitraffic}; CC BY 4.0 & 111 & $3{,}0\unita{s}$ & GPS di bordo, senza quota\\
\bottomrule
\end{tabularx}
\end{table}

\begin{table}[H]\centering\footnotesize
\caption{Copertura della zona dichiarata al livello $0{,}90$ sulla metà di prova, con la divisione per mezzo \prova{P}. Il rumore è simulato dalle schede; le occlusioni sono virtuali.}\label{tab:lp-traiettorie}
\begin{tabularx}{\linewidth}{@{}>{\RaggedRight\arraybackslash}p{3.6cm}>{\RaggedRight\arraybackslash}p{1.6cm}>{\RaggedRight\arraybackslash}p{2.1cm}L>{\RaggedRight\arraybackslash}p{2.0cm}@{}}
\toprule
Classe; rumore simulato & $\tocc$ & Copertura & Bootstrap per mezzo & Wilson, $n$ = mezzi\\
\midrule
aerei, 43 aerei; $\sigma=7{,}5\unita{m}$ a $1\unita{s}$ & 2 / 10 / $30\unita{s}$ & 0,95 / 0,93 / 0,93 & $[0{,}92;0{,}98]$ / $[0{,}90;0{,}96]$ / $[0{,}88;0{,}99]$ & $\ge0{,}83$\\
multirotore, \emph{un} esemplare, 98 voli; radar per droni a $1{,}5\unita{km}$ & 5 / $10\unita{s}$ & 0,99 / 0,96 & per volo: $[0{,}99;1{,}00]$ / $[0{,}95;0{,}97]$ & per volo: $\ge0{,}95$ / $\ge0{,}90$\\
auto, 30 veicoli; radar 4D a $100\unita{m}$ & 2 / 10 / $30\unita{s}$ & 0,93 / 0,87 / 0,86 & $[0{,}90;0{,}95]$ / $[0{,}84;0{,}90]$ / $[0{,}82;0{,}90]$ & $\ge0{,}68$\\
treni, 19 treni; LiDAR a campo stretto a $300\unita{m}$ & 2 / 10 / $30\unita{s}$ & 0,91 / 0,88 / 0,90 & $[0{,}87;0{,}95]$ / $[0{,}84;0{,}92]$ / $[0{,}87;0{,}93]$ & $\ge0{,}67$\\
\bottomrule
\end{tabularx}
\end{table}

La forma più forte che resta vera per gli aerei è: su 43 aerei reali, dopo 2--$30\unita{s}$ senza vederli, la zona dichiarata conteneva l'aereo in circa 9 casi su 10 \prova{P}. Per i droni i dati reali sono due esemplari soltanto, un multirotore e un ala fissa: le cifre valgono per quel multirotore, e per l'ala fissa con guasto indotto la copertura a $10\unita{s}$ scende a $0{,}656$ $[0{,}584;\,0{,}721]$ (21 tratti, divisione per tratto). Nelle virate la zona copre meno: $0{,}74$ per gli aerei e $0{,}72$ per i droni fino a $10\unita{s}$ (divisione per tratto, casi correlati). Per i treni la verità GPS è sporca: con essa la zona tarata diventa larga quanto l'avviso, e non si dichiara una precisione.

\paragraph{L'avviso sulle traiettorie.} Con un inviluppo di $90\unita{m/s}$ per gli aerei l'insieme di possibile presenza non conteneva l'aereo vero in 105 casi su 722 a $2\unita{s}$, e in 167 su 730 a $10\unita{s}$, in 8--11 aerei su 43 \prova{P}: erano decolli e avvicinamenti reali fino a $205\unita{m/s}$. Il principio regge, perché l'avviso non usa l'etichetta; era sbagliato il valore del limite. Per questo l'inviluppo degli aerei in volo è $210\unita{m/s}$; con questo valore la prova sulle traiettorie non è stata rifatta \prova{N}. Anche per i treni l'avviso non conteneva il mezzo in 61 casi su 405 a $2\unita{s}$, in 18 treni su 19, con la verità GPS sporca: la causa non è stata accertata, e per i treni l'avviso non è garantito.

\subsection{Eventi improvvisi e visibilità degradata}\label{sec:lp-eventi}

Gli eventi trattati sono una collisione fra due droni con detriti, l'avaria di un motore con fiammata e fumo durante la corsa di un aereo, la nebbia e la pioggia, il fumo e i detriti in pista, e, come controllo, un passaggio ravvicinato senza collisione. Restano situazioni da osservare e segnalare. Le regole aggiungono alla cella non osservata una \emph{causa} (ostacolo, nube, abbagliamento, meteo, fuori portata, sensore non affidabile, mai guardata), che è un'etichetta in più e non uno stato nuovo. I tre principi non cambiano.

\paragraph{Frammentazione.} Un evento nasce \emph{dedotto} da un innesco di «separazione» (almeno 3 ritorni nuovi confermati entro $2\unita{s}$ attorno a una traccia) o di «contatto» (due tracce più vicine della somma dei mezzi ingombri più $3\sigma$, con almeno un indizio fra ritorni nuovi, residuo fuori porta per 2 passate o una traccia non più rilevata in vista). Diventa \emph{osservato} quando un sensore in salute risolve almeno 3 ritorni distinti per 2 passate; l'avviso è lo stesso nei due casi. Attorno all'ultimo stato comune osservato $p_0$, all'istante $t_a$, i frammenti sono contenuti in un volume $V(\tau)$: in pianta l'unione dei segmenti $[p_0,\,p_0+v_{c,h,i}\tau^*]$ percorsi alle ultime velocità delle tracce d'origine, ingrossata di
\[
r(\tau)=W\tau+(\Delta v+W)\,\tau^*+r_0 ,
\]
con $\Delta v$ la velocità relativa massima dei frammenti, $W$ il vento orizzontale massimo, $\tau^*$ il tempo di volo massimo e $r_0$ l'incertezza iniziale; in quota fra due estremi dati dalla gravità e dal vento verticale.

\begin{proposizione}[Volume dei detriti]\label{prop:lp-detriti}
Per un punto materiale con resistenza isotropa opposta alla velocità relativa all'aria, con coefficiente costante qualunque, vento costante entro $W$ e $W_z$, e velocità iniziale relativa entro $\Delta v$, il frammento sta in $V(\tau)$ per ogni $\tau\le\tau^*$.
\end{proposizione}
\begin{proof}[Traccia della dimostrazione]
Sia $r=v-w$ la velocità relativa all'aria. In orizzontale $\dot r_h=-k|r|\,r_h$: la direzione non cambia e il modulo non cresce, quindi lo spostamento rispetto all'aria è $\lambda\,r_h(0)\,\tau$ con $\lambda\in[0,1]$, e la posizione sta nel segmento ingrossato di $W\tau+(\Delta v+W)\tau$. In verticale, dove $r_z\le0$ vale $\dot r_z\ge-g$ e dove $r_z\ge0$ vale $\dot r_z\le-g$; integrando si hanno i due estremi di quota.
\end{proof}
Il risultato è \prova{D} nel modello per $\tau\le\tau^*$; il taglio a $\tau^*$, il tempo di caduta e il raggio dell'impronta a terra sono sostenuti dal calcolo, non dimostrati. Con $10\,000$ frammenti simulati per caso nessun campione è uscito dal volume (0 su $1\,131\,692$ in 8 casi conformi, 0 su $646\,957$ con raffiche entro i limiti) \prova{S}; il simulatore condivide il modello fisico con la prova, e vede gli errori d'algebra e di codice, non quelli del modello. Se la velocità vera dei frammenti è $1{,}5$ volte quella dichiarata, fra 7 e 88 campioni escono. Restano fuori dal modello la portanza dei pezzi piatti, la corrente ascendente di un incendio, il rotolamento a terra e i frammenti più veloci di $\Delta v$, come la punta di un'elica. Dentro il volume nessuna cella diventa «libera» se non con la probabilità cumulata $0{,}99$ per il frammento minimo ammesso ($5\unita{cm}$); a terra l'impronta resta «zona non verificata: possibili detriti» finché non è ispezionata, e con i sensori della tabella~\ref{tab:lp-sensori} non si scioglie mai da sola.

\paragraph{Nube, abbagliamento, meteo.} Per ogni sensore la nube è l'insieme delle celle fra il primo ritorno anomalo e lo sfondo atteso dalla mappa, stimato su 3 cicli, più $10\unita{m}$; dentro e dietro la nube la cella è non osservata per quel sensore, e solo i radar vedono dietro il fumo (ipotesi prudente). La nube si toglie quando lo sfondo torna visibile per $5\unita{s}$, e le celle tornano libere solo con la probabilità cumulata. L'abbagliamento si riconosce da indizi del sensore stesso (pixel saturi oltre il $5\,\%$ nella camera nel visibile e il $2\,\%$ nella termica, caduta dei ritorni LiDAR oltre il $30\,\%$, fondo radar più alto di $6\unita{dB}$): entro 3 cicli il sensore è dichiarato «affidabilità persa», con la causa, e le sue passate non contano dall'inizio dell'indizio. Con il meteo la portata dichiarata si accorcia per un fattore prudente e non si allunga mai da sola; se il «libero» scende sotto la distanza d'avviso l'uscita dice «avviso non tempestivo», e se poggia su un solo sensore «ridondanza assente». La tabella~\ref{tab:lp-meteo} mostra perché, con nebbia e pioggia forte, tutto poggia sul radar \prova{C}.

\begin{table}[H]\centering\footnotesize
\caption{Distanza del «libero» in metri per sensore e condizione, con i fattori prudenti assunti \prova{C}.}\label{tab:lp-meteo}
\begin{tabularx}{\linewidth}{@{}Lrrrrl@{}}
\toprule
Oggetto, condizione & radar per droni & camera nel visibile & termica & senza radar & avviso a $1490\unita{m}$ senza radar\\
\midrule
multirotore, sereno & 2858 & 2000 & 1200 & 2000 & sì\\
multirotore, pioggia forte & 2000 & 526 & 353 & 526 & no\\
multirotore, nebbia fitta & 2801 & 23 & 15 & 23 & no\\
aereo leggero, nebbia fitta & 4900 & 23 & 15 & 23 & no\\
\bottomrule
\end{tabularx}
\end{table}

\paragraph{Risultati in simulazione.} I conteggi sono per passaggio, con intervalli di Wilson al $95\,\%$; le due serie sono indipendenti nei semi, non nel codice.
\begin{itemize}
\item \emph{Collisione fra droni.} L'evento è dedotto da $1{,}2$ a $2{,}3\unita{s}$ dopo l'urto (mediana $1{,}8\unita{s}$); entro $2\unita{s}$ in 51 passaggi su 60 $[0{,}739;\,0{,}919]$, e con il tipo giusto («contatto») in 50 su 60 $[0{,}720;\,0{,}907]$: in due passaggi è dedotto come «separazione» \prova{S}. La soglia di progetto, entro $2\unita{s}$ in ogni passaggio, non è raggiunta. Nessuna cella è dichiarata «libera» su un frammento: 0 passaggi su 60 $[0;\,0{,}06]$, per costruzione. L'origine dei detriti è presa solo dalle tracce viste negli ultimi $3\unita{s}$: il raggio residuo della zona dei detriti va da $3{,}1$ a $5{,}3\unita{km}$ (mediana $4{,}5\unita{km}$), mai oltre $10\unita{km}$ (0 su 60 $[0;\,0{,}06]$). La zona non si riduce guardando: l'ispezione non ha tolto punti in nessun passaggio (0 su 60), e l'avviso resta su tutte le tracce dentro la zona. Nessun evento è dichiarato prima dell'urto (0 su 60 $[0;\,0{,}06]$), ma la regola delle tracce recenti è stata scritta dopo aver esaminato un passaggio con un evento falso prima dell'urto: questo esito non è una previsione. Con la versione precedente delle regole i frammenti stavano nel volume in una quota di $0{,}9973$, sotto il $0{,}999$ richiesto; con le regole attuali la ripetizione con semi nuovi non è rifatta.
\item \emph{Avaria con fiammata.} L'evento è dedotto in $0{,}2$--$0{,}3\unita{s}$ in 60 passaggi su 60 $[0{,}94;\,1]$, ma solo con un innesco ottimista, perché la saturazione è presa dal punto vero \prova{S}. Le camere sono dichiarate non affidabili entro 3 cicli in 90 passaggi su 90 $[0{,}959;\,1]$. Con le regole attuali le camere sono dichiarate non affidabili entro 3 cicli anche in 30 passaggi su 30. Il mezzo in pista è avvisato a non meno di $140\unita{m}$ dai frammenti veri in 60 passaggi su 60, e con le regole attuali entra davvero nell'impronta in 30 su 30; ma l'impronta di circa $3{,}5\unita{km}$ contiene già il mezzo, e un avviso che parta dal bordo dell'impronta non c'è in nessun passaggio (0 su 60): l'avviso è giusto e non distingue chi è in pericolo: tocca una mediana di 96 tracce.
\item \emph{Nebbia e pioggia.} Nessuna cella «libera» oltre la portata utile, su $4{,}2\cdot10^8$ celle\,$\cdot$\,ciclo, e 0 tracce perse senza avviso su 1080 cicli, con l'inviluppo precedente \prova{S}. Il conto della portata e il sistema usano la stessa formula: la prova non è indipendente.
\item \emph{Fumo e detriti in pista; passaggio ravvicinato.} Nessun falso evento: 0 su 120 e 0 su 60 passaggi col fumo; a $5\unita{m}$ di distanza senza collisione, 0 su 300 $[0;\,0{,}013]$ con l'inviluppo precedente degli aerei e 0 su 30 con quello di $210\unita{m/s}$, anche con le regole attuali (0 su 30 $[0;\,0{,}114]$) \prova{S}.
\end{itemize}
La perdita di un sensore è dichiarata in 460 casi simulati su 460 $[0{,}992;\,1]$, con un ritardo fino a $5{,}3\unita{s}$ (5,3 rivisite del radar del movimento a terra, contro le 3 della soglia di progetto), con l'inviluppo precedente degli aerei \prova{S}; la soglia di 3 rivisite non è quindi rispettata per quel sensore.

\subsection{Prove simulate del caso centrale}\label{sec:lp-prove}

Il sito sorvegliato è descritto da tre scene, ciascuna simulata per passaggio con semi diversi, e con la regola di spegnimento prudente (tabella~\ref{tab:lp-prove}).
\begin{itemize}
\item \emph{Drone dietro un edificio.} Un multirotore a 15--$21\unita{m/s}$ parte a $990\unita{m}$, scende a $12\unita{m}$ di quota e passa per $8\unita{s}$ dietro un edificio di $80\times40\times15\unita{m}$ a $300\unita{m}$ dal sensore. Due configurazioni: radar per droni, radar a pannello, camera nel visibile e identificazione remota; oppure il solo radar per droni.
\item \emph{Aereo in rullaggio dietro un hangar.} Un aereo a $9\unita{m/s}$ sulla via dietro un hangar di $60\times60\times18\unita{m}$, per $17{,}1\unita{s}$; un mezzo a $8{,}3\unita{m/s}$ attraversa la via. Radar del movimento a terra, LiDAR a campo largo e ADS-B, oppure il solo radar.
\item \emph{Drone sotto la copertura del radar d'aeroporto.} Un multirotore scende da $120$ a $30\unita{m}$ sull'asse di una pista; sotto circa $43\unita{m}$ a $2{,}5\unita{km}$ una fascia di alberi ed edifici lo nasconde al radar d'aeroporto, mentre un aereo è in finale.
\end{itemize}

\begin{table}[H]\centering\footnotesize
\caption{Prove simulate del caso centrale \prova{S}. Conteggi per passaggio; intervalli di Wilson al $95\,\%$. «Fuori avviso»: passaggi con almeno un istante in cui il mezzo nascosto è fuori dall'insieme di possibile presenza. Le probabilità di rilevamento sono ipotesi.}\label{tab:lp-prove}
\begin{tabularx}{\linewidth}{@{}>{\RaggedRight\arraybackslash}p{3.0cm}>{\RaggedRight\arraybackslash}p{2.7cm}L@{}}
\toprule
Scena & Configurazione & Esito\\
\midrule
drone dietro un edificio & quattro sensori & avviso tempestivo (servono $806\unita{m}$, traccia confermata a circa $950\unita{m}$) in 200 passaggi su 200 $[0{,}981;\,1]$, e in 100 su 100 $[0{,}963;\,1]$ con semi nuovi e la versione precedente delle regole; fuori avviso 0 su 200, ma solo 22 passaggi su 200 hanno un intervallo nascosto\\
 & solo radar per droni & tempestivo in 200 su 200 $[0{,}981;\,1]$; fuori avviso in 12 su 200 $[0{,}035;\,0{,}102]$ (6 in cui il drone rallenta, 6 in cui si ferma) e in 9 su 100 $[0{,}048;\,0{,}162]$ con semi nuovi e la versione precedente delle regole\\
aereo in rullaggio dietro un hangar & tre sensori & tempestivo in 200 su 200 $[0{,}981;\,1]$; 0 su 200 fuori avviso, 0 senza regione, 0 con «libero» sull'aereo; 34 passaggi su 200 hanno un intervallo nascosto\\
 & solo radar del movimento a terra & tempestivo in 50 su 200 $[0{,}195;\,0{,}314]$, con l'inviluppo precedente: un sensore solo non basta\\
drone sotto la copertura del radar d'aeroporto & radar d'aeroporto e ADS-B & il radar non conferma mai il drone (0 su 200); il drone non è mai dichiarato «libero» (0 su 200 $[0;\,0{,}019]$), e l'uscita dice «zona non verificata per i multirotori». La tempestività non è misurata: lo scenario parte più vicino della distanza richiesta, e non si dà alcuna cifra\\
\bottomrule
\end{tabularx}
\end{table}

Tre precisazioni. L'avviso perde il drone che rallenta o si ferma in circa il $6\,\%$ dei passaggi con il solo radar: la causa non è stata accertata. Il «libero» sull'oggetto con il radar a pannello regge solo con la mappa del sito dichiarata completa (\S\ref{sec:lp-smentiti}): la tempestività di questa tabella non dice che il «libero» sia corretto in ogni caso. Con le regole attuali è rimisurato solo il drone dietro un edificio (200 passaggi per configurazione); il rullaggio e il drone sotto la copertura hanno le cifre della versione precedente.

\paragraph{Altri mezzi, come confronto.} Con i sensori a bordo di un drone (radar 4D, LiDAR a campo largo, camera) l'avviso è tempestivo in 95 passaggi su 100 $[0{,}888;\,0{,}978]$, ma il mezzo esce dall'avviso in 9 su 100, perché il filtro Doppler perde il drone lento. Con l'inviluppo precedente, un'auto a un incrocio dietro un camion è avvisata in tempo in 200 passaggi su 200 $[0{,}981;\,1]$, e un treno in stazione in 186 su 200 $[0{,}886;\,0{,}958]$: il treno che si ferma o torna indietro prima della conferma non è annunciato in tempo \prova{S}.

\paragraph{La scala, riassunta.} In simulazione dalle schede la portata utile passa da $3\unita{m}$ a circa $1\unita{km}$, e la velocità per cui l'avviso è tempestivo da $1{,}44$ a circa $21\unita{m/s}$, per droni, mezzi di piazzale e aerei in rullaggio \prova{S}. Non è una misura su sensori.

\subsection{Aerei in volo: un limite dichiarato}\label{sec:lp-volo}

Con l'inviluppo di $210\unita{m/s}$, l'operatore di sito ha bisogno di $3410\unita{m}$ di preavviso, mentre il «libero» per un multirotore, con passate correlate, arriva a $877\unita{m}$. L'avviso per gli aerei in volo non è tempestivo, e il sistema lo calcola e lo scrive in ogni uscita: «avviso non tempestivo» \prova{C}\prova{S} (caso costruito). Per un treno ad alta velocità il macchinista avrebbe bisogno di $7{,}57\unita{km}$: anche lì l'avviso non è tempestivo, per calcolo \prova{C}. Con un treno sul binario il sistema può avvisare, non garantire l'arresto. Servirebbe un sensore a portata maggiore; nessuna cifra di tempestività è data per gli aerei in volo.

\subsection{Risultati smentiti, come limiti}\label{sec:lp-smentiti}

Le simulazioni hanno trovato casi in cui il sistema non fa ciò che dichiara. Li riportiamo come limiti.
\begin{enumerate}[label=\arabic*.]
\item \textbf{Il «libero» con il radar a pannello dipende dalla mappa.} Nella configurazione del sito, con la regola precedente, una sola passata mancata apriva un «libero» sulla cella del drone in 163 passaggi su 171 $[0{,}910;\,0{,}976]$: il ritorno cadeva nella cella vicina, e legare la finestra alla velocità massima, da solo, non lo toglieva. Con la ritenuta del ritorno per $5\unita{s}$ e con la mappa del sito dichiarata completa, che è un'ipotesi del simulatore, nessun passaggio su 171 ha istanti con «libero» sul drone $[0;\,0{,}022]$ \prova{S}; la ripetizione con semi nuovi non è rifatta. In una configurazione d'aeroporto il «libero» sull'oggetto non compare (0 su 171 $[0;\,0{,}022]$), ma con 8 falsi confermati di mediana per passaggio, e senza avviso tempestivo per gli aerei.
\item \textbf{Mappa dichiarata completa ma sbagliata.} Un edificio che non compare in una mappa dichiarata completa apre un «libero» falso dietro di sé, in 5 casi costruiti su 6 \prova{S}; per questo il prototipo, per impostazione, non dichiara la mappa completa, e lì la cella dietro l'edificio è «non osservata» (6 su 6).
\item \textbf{L'avviso perde il drone che rallenta o si ferma} in 12 passaggi su 200 con il solo radar (\S\ref{sec:lp-prove}); con il filtro Doppler a bordo, in 9 su 100, con la versione precedente delle regole.
\item \textbf{Classi non dichiarate.} Nel sito sorvegliato l'elicottero, l'aereo leggero, l'auto e il pedone non sono né ammessi né esclusi: per loro il «libero» non è definito.
\item \textbf{Eventi improvvisi.} La collisione non è dedotta entro $2\unita{s}$ in ogni passaggio (51 su 60); la zona dei detriti è di alcuni chilometri e non si riduce guardando; l'avviso non distingue chi è in pericolo; la fiammata è rilevata solo con un innesco ottimista (\S\ref{sec:lp-eventi}).
\item \textbf{Inviluppo.} Con $90\unita{m/s}$ l'avviso perdeva aerei reali (8--11 su 43); il valore di $210\unita{m/s}$ non è stato provato di nuovo sulle traiettorie. Il drone veloce oltre i $45\unita{m/s}$ ipotizzati esce dall'avviso.
\item \textbf{Falsi avvisi.} Nelle prove con l'inviluppo precedente si va da $3{,}8$ falsi avvisi all'ora con il solo radar per droni a 1594 all'ora con radar 4D e LiDAR \prova{S}; i falsi ritorni sono un'ipotesi del simulatore.
\end{enumerate}

\subsection{Condizioni di validità a lunga portata}\label{sec:lp-validita}

\begin{riquadro}{Condizioni di validità degli enunciati e delle cifre a lunga portata}
\begin{enumerate}[label=\arabic*.,itemsep=1pt]
\item \textbf{Sensori simulati dalle schede,} con le ipotesi della tabella~\ref{tab:lp-sensori}: probabilità di rilevamento, falsi allarmi, latenze e la sezione radar del multirotore ($0{,}01\unita{m^2}$) non hanno una fonte.
\item \textbf{Sensori fermi, posa nota;} il caso con i sensori a bordo di un drone è un'analisi con posa supposta nota.
\item \textbf{Classi dichiarate dallo scenario} e inviluppo dato: multirotori fino a $45\unita{m/s}$ e $25\unita{m/s^2}$ (ipotesi), aerei in volo fino a $210\unita{m/s}$.
\item \textbf{Mappa degli ostacoli corretta,} oppure sfondo misurato; senza l'una o l'altro non c'è «libero». Le cifre del «libero» sul drone nel sito valgono con la mappa dichiarata completa, che è un'ipotesi.
\item \textbf{Geometrie nominali,} con ostacoli estrusi e oggetti puntiformi; niente riflessioni multiple, uccelli o disturbo strutturato.
\item \textbf{Etichette di classe a mano;} il classificatore è richiesto e supposto esistente, ma non è realizzato: è lasciato ai lavori futuri.
\end{enumerate}
\end{riquadro}
